\documentclass[11pt,a4paper]{article}

\usepackage[T1]{fontenc}
\usepackage[margin=2.4cm]{geometry}
\usepackage{mathpazo}          % Palatino text and maths -- easier to read than Computer Modern
\usepackage{microtype}         % subtler line breaking, fewer ragged edges
\usepackage{amsmath,amssymb}
\usepackage{bm}
\usepackage{booktabs}
\usepackage{array}
\usepackage{xcolor}
\usepackage{titlesec}
\usepackage{parskip}

% --- headings -------------------------------------------------------------
\definecolor{rule}{gray}{0.55}
\titleformat{\section}{\large\bfseries}{}{0pt}{}[{\color{rule}\titlerule}]
\titlespacing*{\section}{0pt}{1.6em}{0.8em}
\titleformat{\subsection}{\bfseries}{}{0pt}{}
\titlespacing*{\subsection}{0pt}{1.2em}{0.4em}

% --- tableaux -------------------------------------------------------------
% Fixed-width columns so the rows of a stacked row operation line up entry by entry.
\newcolumntype{V}{>{\centering\arraybackslash$}p{2.7em}<{$}}
\newcommand{\rowv}[1]{\left[\begin{array}{VVVVVVV|V}#1\end{array}\right]}

% A tableau: caption above, kept on one page, in a slightly smaller size.
\definecolor{pivot}{RGB}{255,236,189}
\newcommand{\piv}[1]{\colorbox{pivot}{$#1$}}   % shades the pivot element
\newenvironment{tableau}[1]
  {\begin{center}\begin{minipage}{\linewidth}\centering\small
   \textbf{#1}\\[0.4em]}
  {\end{minipage}\end{center}}

\title{Engineering Optimisation --- Assignment 1}
\author{}
\date{}

\begin{document}
\maketitle

\section*{Question 1: Robot path around obstacles}

\subsection*{1. Source}

A warehouse robot drives from a pick station to a drop station across a floor
full of pallet stacks, and on the way it must pass a barcode scanner at a fixed
point $\bm{w}$ so the load is logged. Every metre costs time and battery, times
thousands of trips a day. The operator wants the shortest path that passes the
scanner and doesn't hit anything.

\subsection*{2. The three parts}

\textbf{Given data (fixed parameters).} Start $\bm{p}_0$ and goal
$\bm{p}_{N+1}$; the number of waypoints $N$ (chosen beforehand, not a design
variable); $M$ obstacles with centre $\bm{c}_j$ [m] and radius $r_j$ [m]; robot
radius $\rho$ [m]; maximum segment length $\Delta_{\max}$ [m]; scanner position
$\bm{w}$ [m] and the index $q$ of the waypoint that must sit on it; workspace
limits $x_{\min}, x_{\max}, y_{\min}, y_{\max}$ [m]. Start and goal are assumed
to lie clear of every obstacle.

\subsubsection*{(a) Design variables}

\[
  \bm{p}_i = (x_i, y_i) \ [\mathrm{m}], \qquad i = 1,\dots,N
\]

The floor positions of the $N$ waypoints of a polyline path, giving $2N$
design variables in total.

\subsubsection*{(b) Objective (minimise)}

The total path length [m], to be \emph{minimised}:
\[
  f(\bm{p}) = \sum_{i=0}^{N} \lVert \bm{p}_{i+1} - \bm{p}_i \rVert_2
\]

\subsubsection*{(c) Constraints}

\begin{center}
\begin{tabular}{lll}
\toprule
Constraint & Type & Why \\
\midrule
Obstacle clearance     & Behaviour & couples $x_i$ and $y_i$ through a norm \\
Segment length $\le \Delta_{\max}$ & Behaviour & couples two neighbouring waypoints \\
Workspace bounds       & Side      & each limits one variable on its own \\
Scanner gate $\bm{p}_q = \bm{w}$ & Side ($\mathrm{lb}=\mathrm{ub}$) & fixes $x_q$ and $y_q$ one variable at a time \\
\bottomrule
\end{tabular}
\end{center}

The clearance and segment-length constraints work together. Clearance is only
checked at the waypoints, so on its own a segment between two clear waypoints
could still cut through an obstacle. Capping every segment at $\Delta_{\max}$
and inflating the checked radius accordingly closes that gap: if both ends of a
segment of length at most $\Delta_{\max}$ lie at least $R'$ from a centre, every
point of the segment lies at least $\sqrt{R'^2 - \Delta_{\max}^2/4}$ from it.
Choosing
\[
  R_j' = \sqrt{(r_j + \rho)^2 + \Delta_{\max}^2/4}
\]
therefore guarantees that the whole path, not just its waypoints, keeps the
robot off obstacle $j$.

\subsection*{3. Standard form}

\begin{align}
  \min_{\bm{p}} \quad & f(\bm{p}) = \sum_{i=0}^{N} \lVert \bm{p}_{i+1} - \bm{p}_i \rVert_2 \\[0.4em]
  \text{s.t.} \quad
  & g_{ij} = 1 - \frac{\lVert \bm{p}_i - \bm{c}_j \rVert_2}{R_j'} \le 0,
    && i = 1,\dots,N,\ j = 1,\dots,M \\[0.4em]
  & g_i^{\mathrm{seg}} = \frac{\lVert \bm{p}_{i+1} - \bm{p}_i \rVert_2}{\Delta_{\max}} - 1 \le 0,
    && i = 0,\dots,N \\[0.4em]
  & x_{\min} \le x_i \le x_{\max}, \quad y_{\min} \le y_i \le y_{\max},
    && i = 1,\dots,N \\[0.4em]
  & h_1 = x_q - w_x = 0, \qquad h_2 = y_q - w_y = 0
\end{align}

Both inequality families are normalised so that each $g$ is dimensionless and
of order one. In total there are $m = NM + (N+1)$ inequality constraints, $4N$
bounds and $p = 2$ equality constraints.

\subsection*{4. Classification}

\textbf{(a) Nonlinear.} The Euclidean norm appears in the objective and in the
clearance and segment-length constraints.

\textbf{(b) Not convex.} The objective is convex (a sum of norms of affine
functions of $\bm{p}$), the segment-length constraints are convex, and the
bounds and gate are linear. Only obstacle clearance breaks convexity: it asks
each waypoint to stay outside a disc, and a disc exterior isn't convex. Take a
feasible point above the obstacle and one below, and the straight line between
them cuts through. Each way of passing the obstacles (left or right of each one)
gives its own local minimum, so only local optimality can be claimed. In
practice, run the solver from several starting guesses and keep the best.

\newpage
\section*{Question 2: A linear programme by hand}

\subsection*{1. Formulation}

A machine runs 8 h a day for 30 days, so it gives 240 h a month. The hours
available in a group are its number of machines times 240.

\begin{center}
\begin{tabular}{lccccc}
\toprule
 & L & S & D & M & G \\
\midrule
machines available  & 10   & 3   & 4   & 6    & 5    \\
hours per month     & 2400 & 720 & 960 & 1440 & 1200 \\
hours per unit of A & 0.6  & 0.4 & 0.1 & 0.5  & 0.2  \\
hours per unit of B & 0.9  & 0.1 & 0.2 & 0.3  & 0.3  \\
\bottomrule
\end{tabular}
\end{center}

A column is one machine group. A row is one part's recipe.

\textbf{Design variables.} $x_A$ and $x_B$: the parts of type A and type B made
each month [units/month]. Both are continuous. A fractional answer would be
rounded at the end.

\textbf{Objective.} We want $x_A + x_B$ to be as large as possible. This course
always minimises, so we minimise the negative of it: $f = -(x_A + x_B)$. Both
have the same optimum.

\textbf{Constraints.} A group cannot use more hours than it has. Hours used are
(h per unit)$\times$(units per month), which gives h per month, and the limit is
also in h per month.

\begin{align*}
  \min_{x_A,\,x_B} \quad & f = -x_A - x_B && \text{$-$output, units/month} \\
  \text{s.t.} \quad
  & 0.6x_A + 0.9x_B \le 2400 && \text{lathe, h} \\
  & 0.4x_A + 0.1x_B \le 720  && \text{shaper, h} \\
  & 0.1x_A + 0.2x_B \le 960  && \text{drilling, h} \\
  & 0.5x_A + 0.3x_B \le 1440 && \text{milling, h} \\
  & 0.2x_A + 0.3x_B \le 1200 && \text{grinding, h} \\
  & x_A,\ x_B \ge 0
\end{align*}

The five machine limits are \emph{behaviour} constraints. They limit how the
shop performs, and each one uses both variables. The two sign conditions are
\emph{side} constraints: each limits a single variable, so they are bounds.

\subsection*{2. Standard form with slack variables}

One slack per constraint turns every $\le$ into an $=$:

\begin{align*}
  & 0.6x_A + 0.9x_B + s_L = 2400 \\
  & 0.4x_A + 0.1x_B + s_S = 720 \\
  & 0.1x_A + 0.2x_B + s_D = 960 \\
  & 0.5x_A + 0.3x_B + s_M = 1440 \\
  & 0.2x_A + 0.3x_B + s_G = 1200 \\
  & x_A,\ x_B,\ s_L,\ s_S,\ s_D,\ s_M,\ s_G \ge 0
\end{align*}

\textbf{What each slack measures.} A slack is the time that group stands idle
each month [h/month]: $s_L$ unused lathe hours, $s_S$ shaper, $s_D$ drilling,
$s_M$ milling, $s_G$ grinding. If a slack is zero, the group is fully booked and
its constraint is active. If it is positive, the group has that many hours to
spare.

\subsection*{3. Simplex tableau [Bonus]}

Move the objective across: $-x_A - x_B + 0s_L + 0s_S + 0s_D + 0s_M + 0s_G + (-f) = 0$.
The rows are $R_1,\dots,R_5$ (L, S, D, M, G) and $R_{-f}$. The pivot element is
shaded in each tableau.

\begin{tableau}{Initial tableau}
\begin{tabular}{l|rrrrrrr|r}
basis & $x_A$ & $x_B$ & $s_L$ & $s_S$ & $s_D$ & $s_M$ & $s_G$ & RHS \\
\hline
$s_L$ & 0.6 & 0.9 & 1 & 0 & 0 & 0 & 0 & 2400 \\
$s_S$ & \piv{0.4} & 0.1 & 0 & 1 & 0 & 0 & 0 & 720 \\
$s_D$ & 0.1 & 0.2 & 0 & 0 & 1 & 0 & 0 & 960 \\
$s_M$ & 0.5 & 0.3 & 0 & 0 & 0 & 1 & 0 & 1440 \\
$s_G$ & 0.2 & 0.3 & 0 & 0 & 0 & 0 & 1 & 1200 \\
\hline
$-f$  & $-1$ & $-1$ & 0 & 0 & 0 & 0 & 0 & 0 \\
\end{tabular}
\end{tableau}

The basis is $\{s_L, s_S, s_D, s_M, s_G\}$, so $x_A = x_B = 0$ and every slack
holds its full right-hand side. We make nothing and every machine stands idle.
That is a real vertex, and a bad plan.

\textbf{Iteration 1.} The most negative entry in the $-f$ row is $-1$, and
$x_A$ and $x_B$ tie, because each part adds one unit of output. Either may
enter. We choose $x_A$. The ratio test on the $x_A$ column:

\begin{center}
\begin{tabular}{lll}
\toprule
Row & Ratio & Meaning \\
\midrule
$s_L$ (lathe)    & $2400/0.6 = 4000$ & the lathe would last for 4000 of part A \\
$s_S$ (shaper)   & $720/0.4 = 1800$  & smallest --- the shaper runs out first \\
$s_D$ (drilling) & $960/0.1 = 9600$  & drilling would last for 9600 \\
$s_M$ (milling)  & $1440/0.5 = 2880$ & milling would last for 2880 \\
$s_G$ (grinding) & $1200/0.2 = 6000$ & grinding would last for 6000 \\
\bottomrule
\end{tabular}
\end{center}

So $s_S$ leaves and the pivot element is the 0.4.

\emph{Step A} --- divide the pivot row by 0.4:
\[
  R_2 \leftarrow R_2/0.4 \;=\; \rowv{1 & 0.25 & 0 & 2.5 & 0 & 0 & 0 & 1800}
\]

\emph{Step B} --- clear the $x_A$ column with that new $R_2$. The first one in
full; the others work the same way:
\begin{align*}
  R_1 \leftarrow R_1 - 0.6R_2: \quad
  & \rowv{0.6 & 0.9 & 1 & 0 & 0 & 0 & 0 & 2400} \\
  -\,0.6\times & \rowv{1 & 0.25 & 0 & 2.5 & 0 & 0 & 0 & 1800} \\
  = & \rowv{0 & 0.75 & 1 & -1.5 & 0 & 0 & 0 & 1320} \\
  R_3 \leftarrow R_3 - 0.1R_2 \;=\; & \rowv{0 & 0.175 & 0 & -0.25 & 1 & 0 & 0 & 780} \\
  R_4 \leftarrow R_4 - 0.5R_2 \;=\; & \rowv{0 & 0.175 & 0 & -1.25 & 0 & 1 & 0 & 540} \\
  R_5 \leftarrow R_5 - 0.2R_2 \;=\; & \rowv{0 & 0.25 & 0 & -0.5 & 0 & 0 & 1 & 840} \\
  R_{-f} \leftarrow R_{-f} + R_2 \;=\; & \rowv{0 & -0.75 & 0 & 2.5 & 0 & 0 & 0 & 1800}
\end{align*}

\begin{tableau}{Tableau after iteration 1}
\begin{tabular}{l|rrrrrrr|r}
basis & $x_A$ & $x_B$ & $s_L$ & $s_S$ & $s_D$ & $s_M$ & $s_G$ & RHS \\
\hline
$s_L$ & 0 & \piv{0.75} & 1 & $-1.5$  & 0 & 0 & 0 & 1320 \\
$x_A$ & 1 & 0.25  & 0 & 2.5     & 0 & 0 & 0 & 1800 \\
$s_D$ & 0 & 0.175 & 0 & $-0.25$ & 1 & 0 & 0 & 780 \\
$s_M$ & 0 & 0.175 & 0 & $-1.25$ & 0 & 1 & 0 & 540 \\
$s_G$ & 0 & 0.25  & 0 & $-0.5$  & 0 & 0 & 1 & 840 \\
\hline
$-f$  & 0 & $-0.75$ & 0 & 2.5 & 0 & 0 & 0 & 1800 \\
\end{tabular}
\end{tableau}

We are now at the vertex $(x_A, x_B) = (1800, 0)$: make only part A, 1800 a
month. $s_S = 0$, so the shaper constraint is now active. The $-f$ row still has
a negative entry, $-0.75$ under $x_B$, so making some part B would raise output.

\textbf{Iteration 2.} $x_B$ enters. The ratio test on the $x_B$ column:

\begin{center}
\begin{tabular}{lll}
\toprule
Row & Ratio & Meaning \\
\midrule
$s_L$ (lathe)    & $1320/0.75 = 1760$       & smallest --- the lathe runs out first \\
$x_A$            & $1800/0.25 = 7200$       & we would run out of part A to trade away \\
$s_D$ (drilling) & $780/0.175 \approx 4457$ & drilling still has room \\
$s_M$ (milling)  & $540/0.175 \approx 3086$ & milling still has room \\
$s_G$ (grinding) & $840/0.25 = 3360$        & grinding still has room \\
\bottomrule
\end{tabular}
\end{center}

So $s_L$ leaves and the pivot element is the 0.75.

\emph{Step A} --- divide the pivot row by 0.75:
\[
  R_1 \leftarrow R_1/0.75 \;=\; \rowv{0 & 1 & \tfrac{4}{3} & -2 & 0 & 0 & 0 & 1760}
\]

\emph{Step B} --- clear the $x_B$ column with that new $R_1$:
\begin{align*}
  R_2 \leftarrow R_2 - 0.25R_1: \quad
  & \rowv{1 & 0.25 & 0 & 2.5 & 0 & 0 & 0 & 1800} \\
  -\,0.25\times & \rowv{0 & 1 & \tfrac{4}{3} & -2 & 0 & 0 & 0 & 1760} \\
  = & \rowv{1 & 0 & -\tfrac{1}{3} & 3 & 0 & 0 & 0 & 1360} \\
  R_3 \leftarrow R_3 - 0.175R_1 \;=\; & \rowv{0 & 0 & -\tfrac{7}{30} & 0.1 & 1 & 0 & 0 & 472} \\
  R_4 \leftarrow R_4 - 0.175R_1 \;=\; & \rowv{0 & 0 & -\tfrac{7}{30} & -0.9 & 0 & 1 & 0 & 232} \\
  R_5 \leftarrow R_5 - 0.25R_1 \;=\; & \rowv{0 & 0 & -\tfrac{1}{3} & 0 & 0 & 0 & 1 & 400} \\
  R_{-f} \leftarrow R_{-f} + 0.75R_1 \;=\; & \rowv{0 & 0 & 1 & 1 & 0 & 0 & 0 & 3120}
\end{align*}

\begin{tableau}{Final tableau}
\renewcommand{\arraystretch}{1.4}
\begin{tabular}{l|rrrrrrr|r}
basis & $x_A$ & $x_B$ & $s_L$ & $s_S$ & $s_D$ & $s_M$ & $s_G$ & RHS \\
\hline
$x_B$ & 0 & 1 & $\tfrac{4}{3}$   & $-2$   & 0 & 0 & 0 & 1760 \\
$x_A$ & 1 & 0 & $-\tfrac{1}{3}$  & 3      & 0 & 0 & 0 & 1360 \\
$s_D$ & 0 & 0 & $-\tfrac{7}{30}$ & 0.1    & 1 & 0 & 0 & 472 \\
$s_M$ & 0 & 0 & $-\tfrac{7}{30}$ & $-0.9$ & 0 & 1 & 0 & 232 \\
$s_G$ & 0 & 0 & $-\tfrac{1}{3}$  & 0      & 0 & 0 & 1 & 400 \\
\hline
$-f$  & 0 & 0 & 1 & 1 & 0 & 0 & 0 & 3120 \\
\end{tabular}
\end{tableau}

No entry in the $-f$ row is negative, so no neighbouring vertex is better. An LP
is convex, so no vertex anywhere is better either. The optimum is
\[
  x_A^\star = 1360, \qquad x_B^\star = 1760 \ \text{units/month}, \qquad
  f^\star = -3120,
\]
a total of 3120 parts a month. The answer is whole numbers, so no rounding is
needed. Putting it back into the original constraints: the lathe uses
$0.6(1360) + 0.9(1760) = 2400$ h and the shaper $0.4(1360) + 0.1(1760) = 720$ h.
Both are exactly at their limit.

% \subsection*{4. Shadow prices [Bonus]}

% The shadow prices are the slack columns of the final $-f$ row:

% \begin{center}
% \begin{tabular}{lcccl}
% \toprule
% Constraint & $\lambda$ [units/h] & Status & Slack [h] & What it means \\
% \midrule
% lathe    & 1 & active   & 0   & bottleneck \\
% shaper   & 1 & active   & 0   & bottleneck \\
% drilling & 0 & inactive & 472 & 472 of 960 h idle \\
% milling  & 0 & inactive & 232 & 232 of 1440 h idle \\
% grinding & 0 & inactive & 400 & 400 of 1200 h idle \\
% \bottomrule
% \end{tabular}
% \end{center}

% \textbf{The bottlenecks} are the lathe and the shaper. Both slacks are zero, so
% all 2400 lathe hours and all 720 shaper hours are used. One more hour on either
% gives one more part a month.

% \textbf{Which machine to buy.} A machine adds 240 h a month, so a lathe or a
% shaper would each add $1 \times 240 = 240$ parts a month. That holds for the
% whole machine, not only the first hour: with 240 extra hours in either group,
% every basic variable in the final tableau stays positive, so the basis and the
% prices do not change. The tightest case is milling, which falls to 16 h idle
% when a shaper is added. The two are worth the same, so the choice comes down to
% price. I would not buy a drill, a mill or a grinder.

% \textbf{Why three prices are zero.} Drilling, milling and grinding already stand
% idle for 472, 232 and 400 hours a month. A 473rd idle drilling hour is worth
% nothing, because drilling time is not what holds the plan back. An inactive
% constraint prices at zero, and that zero says where not to spend money.

\newpage
\section*{Question 3: A larger problem, solved in MATLAB}

\subsection*{1. Formulation}

The decision is which farm supplies which vegetable, so there is one variable
per farm--vegetable pair:
\[
  x_{fv} = \text{tonnes of vegetable } v \text{ bought from farm } f,
  \qquad f = 1,\dots,4, \quad v = 1,\dots,5,
\]
which is $4 \times 5 = 20$ design variables [t/month]. With $c_{fv}$ the price
in \$/t, the problem is
\begin{align*}
  \min \quad & f = \sum_f \sum_v c_{fv} x_{fv} && \text{total cost, \$} \\
  \text{s.t.} \quad
  & \sum_v x_{fv} \le \text{supply}_f && \text{4 rows, one per farm} \\
  & \sum_f x_{fv} \ge \text{demand}_v && \text{5 rows, one per vegetable} \\
  & x_{fv} \ge 0
\end{align*}

The supply rows are $\le$ because a farm cannot deliver more than its maximum.
The demand rows are $\ge$ because the store must buy at least what it needs.
For \texttt{linprog} the demand rows are multiplied by $-1$ to make them $\le$.

\subsection*{2. Solution}

The MATLAB file is \texttt{A1\_q3.m}. It builds $c$, $A$ and $b$ from the price
table instead of typing out the coefficients. The buying plan [t]:

\begin{center}
\begin{tabular}{lrrrrrr}
\toprule
 & Potato & Tomato & Okra & Eggplant & Spinach & Total \\
\midrule
Farm 1 & 40  & 0  & 0  & 0  & 0  & 40 \\
Farm 2 & 0   & 0  & 20 & 0  & 20 & 40 \\
Farm 3 & 0   & 0  & 0  & 80 & 20 & 100 \\
Farm 4 & 60  & 60 & 0  & 0  & 0  & 120 \\
\midrule
Total  & 100 & 60 & 20 & 80 & 40 & \\
\bottomrule
\end{tabular}
\end{center}

Total cost: \$173\,000 a month. Every vegetable is bought in exactly the
required amount, since buying more only costs money.

\subsection*{3. Which farms are at capacity}

Farm 3 (100 t of 100) and farm 4 (120 t of 120) are at capacity. Farms 1 and 2
are not: they deliver 40 t each, out of 180 t and 200 t.

Potato and spinach are each bought from two farms. Potato is cheapest at farm 4
(\$150/t), but farm 4 is full after 60 t of potato and 60 t of tomato, so the
remaining 40 t comes from the next cheapest farm, farm 1 (\$200/t). Spinach is
cheapest at farm 3 (\$1000/t), but farm 3 is full after 80 t of eggplant and
20 t of spinach, so the other 20 t comes from farm 2 (\$1100/t). So a vegetable
is split only when the cheap farm runs out of capacity.

\end{document}
